% Part: first-order-logic
% Chapter: sequent-calculus
% Section: proving-things

\documentclass[../../../include/open-logic-section]{subfiles}

\begin{document}

\iftag{FOL}
      {\olfileid{fol}{seq}{pro}}
      {\olfileid{pl}{seq}{pro}}

\olsection{\usetoken{P}{derivation}નાં દાખલા}

\begin{ex}
સિક્વન્ટ $!A \land !B \Sequent !A$ની $\Log{LK}$-!!{derivation} આપો.

પહેલાં ઇચ્છિત અંતિમ સિક્વન્ટને !!{derivation}ની નીચે લખીએ.
\begin{prooftree}
\AxiomC{}
\UnaryInf$!A\land !B \fCenter !A$
\end{prooftree}
હવે કયા પ્રકારના અનુમાનનું નીચલું સિક્વન્ટ આ સ્વરૂપનું હોઈ શકે તે શોધવું
પડશે. તે કોઈ સંરચનાત્મક નિયમ હોઈ શકે, પરંતુ પહેલાં તાર્કિક નિયમ શોધવો
સારો ઉપાય છે. નીચલા સિક્વન્ટમાં આવતો એકમાત્ર તાર્કિક સંયોજક $\land$ છે;
તેથી આપણે $\land$નો નિયમ શોધીએ છીએ. $\land$નું ચિહ્ન પૂર્વાંગમાં આવે છે,
તેથી આપણે \LeftR{\land} નિયમ જોઈએ છીએ.
\begin{prooftree}
\AxiomC{}
\RightLabel{\LeftR{\land}}
\UnaryInf$!A\land !B \fCenter !A$
\end{prooftree}
\LeftR{\land} અનુમાનનું ઉપરનું સિક્વન્ટ કયું હતું તેની બે શક્યતાઓ છે:
$!A \Sequent !A$ અથવા $!B \Sequent !A$. સ્પષ્ટ છે કે
$!A \Sequent !A$ આરંભિક સિક્વન્ટ છે, જ્યારે સામાન્ય રીતે
$!B \Sequent !A$ નિષ્પન્ન થઈ શકતું નથી. તેથી ઉપરનું સિક્વન્ટ ભરીએ:
\begin{prooftree}
\Axiom$!A \fCenter !A$
\RightLabel{\LeftR{\land}}
\UnaryInf$!A\land !B \fCenter !A$
\end{prooftree}
હવે આપણી પાસે સિક્વન્ટ $!A \land !B \Sequent !A$ની યોગ્ય
$\Log{LK}$-!!{derivation} છે.
\end{ex}

\begin{ex}
સિક્વન્ટ $\lnot !A \lor !B \Sequent !A \lif !B$ની
$\Log{LK}$-!!{derivation} આપો.

ઇચ્છિત અંતિમ સિક્વન્ટને !!{derivation}ની નીચે લખીને શરૂ કરીએ.
\begin{prooftree}
\AxiomC{}
\UnaryInf$\lnot !A \lor !B \fCenter !A \lif !B$
\end{prooftree}
આ અંતિમ સિક્વન્ટ આપતો હોય એવો તાર્કિક નિયમ શોધવા માટે તેમાંના તાર્કિક
સંયોજકો જોઈએ: $\lnot$, $\lor$ અને $\lif$. અત્યારે આપણે માત્ર $\lor$ અને
$\lif$ની જ ચિંતા કરીએ છીએ, કારણ કે તેઓ અંતિમ સિક્વન્ટમાંનાં
!!{sentence}sના !!{main operator}s છે. $\lnot$ બીજા સંયોજકના વ્યાપમાં
છે, તેથી તેને પછી સંભાળીશું. અંતિમ અનુમાન માટેના તાર્કિક નિયમોની આપણી
શક્યતાઓ \LeftR{\lor} અને \RightR{\lif} છે. બંનેમાંથી ગમે તે નિયમ લઈ
શકાય, પરંતુ આપણે \RightR{\lif} લઈએ છીએ (બીજું કોઈ કારણ ન હોય તોપણ તે
વૃક્ષને બે શાખામાં વહેંચવાનું થોડું મોડું કરવા દે છે). નીચલું સિક્વન્ટ
આપતા \RightR{\lif} અનુમાનના સ્વરૂપ મુજબ તે આવું જ હોવું જોઈએ:
\begin{prooftree}
\AxiomC{}
\UnaryInf$ !A, \lnot !A \lor !B \fCenter !B $
\RightLabel{\RightR{\lif}} \UnaryInf$ \lnot !A \lor !B \fCenter !A \lif !B $
\end{prooftree}

પૂર્વાંગમાં $\lnot !A \lor !B$ને બહારના છેડે ખસેડીએ તો
\LeftR{\lor} નિયમ લગાડી શકીએ. તેના પ્રરૂપ મુજબ તે નીચે પ્રમાણે બે ઉપરના
સિક્વન્ટોમાં વહેંચાવું જોઈએ:
\begin{prooftree}
\AxiomC{}
\UnaryInf$\lnot !A, !A \fCenter !B$
\AxiomC{}
\UnaryInf$!B, !A \fCenter !B$
\RightLabel{\LeftR{\lor}}
\BinaryInf$ \lnot !A \lor !B, !A \fCenter !B $
\RightLabel{\LeftR{\Exchange}}
\UnaryInf$ !A, \lnot !A \lor !B \fCenter !B $
\RightLabel{\RightR{\lif}}
\UnaryInf$ \lnot !A \lor !B \fCenter !A \lif !B $
\end{prooftree}
\sourcecorrection{OLSEQ-001}{મૂળની \texttt{proving-things.tex}માં આ
ઉદાહરણનાં ચાર વૃક્ષોમાં પૂર્વાંગના બે !!{sentence}sની અદલાબદલી કરતી
રેખાનું નામ $\RightR{\Exchange}$ આપ્યું છે. નિયમના સ્વરૂપ અને સાથના
વર્ણન મુજબ તેને ચારેય સ્થળે $\LeftR{\Exchange}$ કર્યું છે.}
યાદ રાખો કે આપણે પાછળની દિશામાં આરંભિક સિક્વન્ટો સુધી પહોંચવાનો પ્રયત્ન
કરી રહ્યા છીએ; હવે આપણે તેની ઘણી નજીક દેખાઈએ છીએ!{} જમણી શાખા આરંભિક
સિક્વન્ટથી માત્ર એક શિથિલીકરણ અને એક અદલાબદલી દૂર છે; પછી તે પૂરી થાય છે:
\begin{prooftree}
\AxiomC{}
\UnaryInf$\lnot !A, !A \fCenter !B$
\Axiom$!B \fCenter !B$
\RightLabel{\LeftR{\Weakening}}
\UnaryInf$!A, !B \fCenter !B$
\RightLabel{\LeftR{\Exchange}}
\UnaryInf$!B, !A \fCenter !B$
\RightLabel{\LeftR{\lor}}
\BinaryInf$\lnot !A \lor !B, !A \fCenter !B $
\RightLabel{\LeftR{\Exchange}}
\UnaryInf$ !A, \lnot !A \lor !B \fCenter !B $
\RightLabel{\RightR{\lif}}
\UnaryInf$ \lnot !A \lor !B \fCenter !A \lif !B $
\end{prooftree}

હવે ડાબી શાખા જોઈએ. તેના કોઈ !!{sentence}માંનો એકમાત્ર તાર્કિક
સંયોજક પૂર્વાંગનાં !!{sentence}sમાં આવતું $\lnot$ ચિહ્ન છે, તેથી આપણે
\LeftR{\lnot} નિયમનો કોઈ ઉપયોગ શોધીએ છીએ.
\begin{prooftree}
\AxiomC{}
\UnaryInf$ !A \fCenter !B, !A$
\RightLabel{\LeftR{\lnot}}
\UnaryInf$\lnot !A, !A \fCenter !B$
\Axiom$!B \fCenter !B$
\RightLabel{\LeftR{\Weakening}}
\UnaryInf$!A, !B \fCenter !B$
\RightLabel{\LeftR{\Exchange}}
\UnaryInf$!B, !A \fCenter !B$
\RightLabel{\LeftR{\lor}}
\BinaryInf$\lnot !A \lor !B, !A \fCenter !B $
\RightLabel{\LeftR{\Exchange}}
\UnaryInf$ !A, \lnot !A \lor !B \fCenter !B $
\RightLabel{\RightR{\lif}}
\UnaryInf$ \lnot !A \lor !B \fCenter !A \lif !B $
\end{prooftree}
જેમ જમણી શાખા પૂરી કરી હતી તેમ આ ડાબી શાખા પણ પૂરી થવાથી માત્ર એક
શિથિલીકરણ અને એક અદલાબદલી દૂર છે.
\begin{prooftree}
\Axiom$!A \fCenter !A$
\RightLabel{\RightR{\Weakening}}
\UnaryInf$ !A \fCenter !A, !B$
\RightLabel{\RightR{\Exchange}}
\UnaryInf$ !A \fCenter !B, !A$
\RightLabel{\LeftR{\lnot}}
\UnaryInf$\lnot !A, !A \fCenter !B$
\Axiom$!B \fCenter !B$
\RightLabel{\LeftR{\Weakening}}
\UnaryInf$!A, !B \fCenter !B$
\RightLabel{\LeftR{\Exchange}}
\UnaryInf$!B, !A \fCenter !B$
\RightLabel{\LeftR{\lor}}
\BinaryInf$\lnot !A \lor !B, !A \fCenter !B $
\RightLabel{\LeftR{\Exchange}}
\UnaryInf$ !A, \lnot !A \lor !B \fCenter !B $
\RightLabel{\RightR{\lif}}
\UnaryInf$ \lnot !A \lor !B \fCenter !A \lif !B $
\end{prooftree}
\end{ex}

\begin{ex}
સિક્વન્ટ $\lnot !A \lor \lnot !B \Sequent \lnot (!A \land !B)$ની
$\Log{LK}$-!!{derivation} આપો.

ઉપરની રીતો વાપરીને ઇચ્છિત અંતિમ સિક્વન્ટને નીચે લખીને શરૂ કરીએ.
\begin{prooftree}
\AxiomC{}
\UnaryInf$ \lnot !A \lor \lnot !B \fCenter \lnot (!A \land !B) $
\end{prooftree}
અંતિમ સિક્વન્ટનાં !!{sentence}sમાં ઉપલબ્ધ મુખ્ય સંયોજકો $\lor$ અને
$\lnot$ છે. અહીં \LeftR{\lor} અથવા \RightR{\lnot}, બંનેમાંથી ગમે તે
નિયમ લગાડી શકાય; પરંતુ થોડી વાર માટે બે શાખામાં વહેંચવાનું ટાળવા આપણે
\RightR{\lnot}થી શરૂ કરીએ છીએ:
\begin{prooftree}
\AxiomC{}
\UnaryInf$!A \land !B, \lnot !A \lor \lnot !B \fCenter $
\RightLabel{\RightR{\lnot}}
\UnaryInf$\lnot !A \lor \lnot !B \fCenter \lnot (!A \land !B)$
\end{prooftree}
હવે \LeftR{\land} કે \LeftR{\lor}, કયો નિયમ જોવો તેની પસંદગી છે.
\LeftR{\land} નિયમ લગાડીએ તો શું થાય તે જોઈએ: આપણે કાં તો સિક્વન્ટ
$!A, \lnot !A \lor \lnot !B \Sequent \quad$થી અથવા સિક્વન્ટ
$!B, \lnot !A \lor \lnot !B \Sequent \quad$થી શરૂ કરી શકીએ છીએ.
!!{derivation} એ $!A$ અને $!B$ની દૃષ્ટિએ સમમિત છે, તેથી પહેલું લઈએ:
\sourcecorrection{OLSEQ-002}{મૂળની \texttt{proving-things.tex}માં આ
વાક્યનાં બંને વૈકલ્પિક સિક્વન્ટોમાં વિયોજનનો બીજો ઘટક $!B$ લખાયો છે.
ઇચ્છિત અંતિમ સિક્વન્ટ અને પછીનાં બંને શાખા-વૃક્ષો મુજબ તે બંને સ્થળે
$\lnot !B$ હોવો જોઈએ.}
\begin{prooftree}
\AxiomC{}
\UnaryInf$!A, \lnot !A \lor \lnot !B \fCenter $
\RightLabel{\LeftR{\land}}
\UnaryInf$!A \land !B, \lnot !A \lor \lnot !B \fCenter $
\RightLabel{\RightR{\lnot}}
\UnaryInf$\lnot !A \lor \lnot !B \fCenter \lnot (!A \land !B)$
\end{prooftree}
!!{derivation} ભરવાનું આગળ વધારતાં આપણને સમસ્યા મળે છે:
\begin{prooftree}
\Axiom$!A \fCenter !A$
\RightLabel{\LeftR{\lnot}}
\UnaryInf$ \lnot !A, !A \fCenter$
\AxiomC{}
\RightLabel{?}
\UnaryInf$!A \fCenter !B$
\RightLabel{\LeftR{\lnot}}
\UnaryInf$ \lnot !B, !A \fCenter$
\RightLabel{\LeftR{\lor}}
\BinaryInf$\lnot !A \lor \lnot !B, !A \fCenter $
\RightLabel{\LeftR{\Exchange}}
\UnaryInf$!A, \lnot !A \lor \lnot !B \fCenter $
\RightLabel{\LeftR{\land}}
\UnaryInf$!A \land !B, \lnot !A \lor \lnot !B \fCenter $
\RightLabel{\RightR{\lnot}}
\UnaryInf$\lnot !A \lor \lnot !B \fCenter \lnot (!A \land !B)$
\end{prooftree}
જમણી શાખાના ઉપરના ભાગને હવે વધુ ઘટાડવો શક્ય નથી અને સંરચનાત્મક
અનુમાનો વડે તેને આરંભિક સિક્વન્ટ સુધી લઈ જઈ શકાતો નથી. તેથી આ સાચો
માર્ગ નથી. એટલે ઉપર \LeftR{\land} લગાડવું ભૂલ હતું. એ પહેલાં જે હતું
ત્યાં પાછાં જઈને તેના બદલે \LeftR{\lor} નિયમ લગાડીએ તો મળે છે:
\begin{prooftree}
\AxiomC{}
\UnaryInf$\lnot !A, !A \land !B \fCenter $

\AxiomC{}
\UnaryInf$\lnot !B, !A \land !B \fCenter $

\RightLabel{\LeftR{\lor}}
\BinaryInf$\lnot !A \lor \lnot !B, !A \land !B \fCenter $
\RightLabel{\LeftR{\Exchange}}
\UnaryInf$!A \land !B, \lnot !A \lor \lnot !B \fCenter $
\RightLabel{\RightR{\lnot}}
\UnaryInf$\lnot !A \lor \lnot !B \fCenter \lnot (!A \land !B)$
\end{prooftree}
આગળની જેમ દરેક શાખા પૂરી કરતાં મળે છે:
\begin{prooftree}
\Axiom$ !A \fCenter!A$
\RightLabel{\LeftR{\land}}
\UnaryInf$!A \land !B \fCenter !A$
\RightLabel{\LeftR{\lnot}}
\UnaryInf$\lnot !A, !A \land !B \fCenter $

\Axiom$ !B \fCenter !B$
\RightLabel{\LeftR{\land}}
\UnaryInf$!A \land !B \fCenter !B$
\RightLabel{\LeftR{\lnot}}
\UnaryInf$\lnot !B, !A \land !B \fCenter $

\RightLabel{\LeftR{\lor}}
\BinaryInf$\lnot !A \lor \lnot !B, !A \land !B \fCenter $
\RightLabel{\LeftR{\Exchange}}
\UnaryInf$!A \land !B, \lnot !A \lor \lnot !B \fCenter $
\RightLabel{\RightR{\lnot}}
\UnaryInf$\lnot !A \lor \lnot !B \fCenter \lnot (!A \land !B)$
\end{prooftree}
(આ પગલાંમાં $\land$ના નિયમો $\lnot$ના નિયમોથી નીચે લગાડ્યા હોત તોપણ
યોગ્ય !!{derivation} મળી હોત.)
\end{ex}

\begin{ex}
હજી સુધી આપણે સંકોચનનો નિયમ વાપર્યો નથી, પરંતુ ક્યારેક તે જરૂરી બને છે.
તેનું એક ઉદાહરણ જોઈએ. માનો કે આપણે $\quad \Sequent !A \lor \lnot !A$
સાબિત કરવા માગીએ છીએ. $\RightR{\lor}$ને ઊલટી દિશામાં લગાડવાથી નીચેની
બે !!{derivation}sમાંથી કોઈ એક મળે:
\begin{prooftree}
\AxiomC{}
\UnaryInf$ \fCenter !A$
\RightLabel{\RightR{\lor}}
\UnaryInf$ \fCenter !A \lor \lnot !A$
\DisplayProof\qquad\bottomAlignProof
\AxiomC{}
\UnaryInf$!A \fCenter $
\RightLabel{\RightR{\lnot}}
\UnaryInf$ \fCenter \lnot !A$
\RightLabel{\RightR{\lor}}
\UnaryInf$ \fCenter !A \lor \lnot !A$
\end{prooftree}
આ બંનેમાંથી એકેયનું ઉપરનું છેડું આરંભિક સિક્વન્ટ નથી. યુક્તિ એ સમજવાની
છે કે સંકોચનનો નિયમ !!a{sentence}ની બે નકલોને એકમાં ભેગી કરવા દે છે—અને
સાબિતી શોધતી વખતે, એટલે કે નીચેથી ઉપર જતી વખતે, આપણે આધાર-સિક્વન્ટમાં
$!A \lor \lnot !A$ની એક નકલ રાખી શકીએ છીએ; દાખલા તરીકે:
\begin{prooftree}
\AxiomC{}
\UnaryInf$ \fCenter !A \lor \lnot !A, !A$
\RightLabel{\RightR{\lor}}
\UnaryInf$ \fCenter !A \lor \lnot !A, !A \lor \lnot !A$
\RightLabel{\RightR{\Contraction}}
\UnaryInf$ \fCenter !A \lor \lnot !A$
\end{prooftree}
હવે $\RightR{\lor}$ને બીજી વાર લગાડી શકાય અને~$\lnot !A$ પણ મળે છે;
તેનાથી પૂર્ણ !!{derivation} મળે છે.
\begin{prooftree}
\Axiom$!A \fCenter !A$
\RightLabel{\RightR{\lnot}}
\UnaryInf$\fCenter !A, \lnot !A$
\RightLabel{\RightR{\lor}}
\UnaryInf$\fCenter !A, !A \lor \lnot !A$
\RightLabel{\RightR{\Exchange}}
\UnaryInf$ \fCenter !A \lor \lnot !A, !A$
\RightLabel{\RightR{\lor}}
\UnaryInf$ \fCenter !A \lor \lnot !A, !A \lor \lnot !A$
\RightLabel{\RightR{\Contraction}}
\UnaryInf$ \fCenter !A \lor \lnot !A$
\end{prooftree}
\end{ex}

\begin{prob}
નીચેના સિક્વન્ટોની !!{derivation}s આપો:
\begin{enumerate}
\item $!A \land (!B \land !C) \Sequent (!A \land !B) \land !C$.
\item $!A \lor (!B \lor !C) \Sequent (!A \lor !B) \lor !C$.
\item $!A \lif (!B \lif !C) \Sequent !B \lif (!A \lif !C)$.
\item $!A \Sequent \lnot\lnot !A$.
\end{enumerate}
\end{prob}

\begin{prob}
નીચેના સિક્વન્ટોની !!{derivation}s આપો:
\begin{enumerate}
\item $(!A \lor !B) \lif !C \Sequent !A \lif !C$.
\item $(!A \lif !C) \land (!B \lif !C) \Sequent (!A \lor !B) \lif !C$.
\item $\Sequent \lnot(!A \land \lnot !A)$.
\item $!B \lif !A \Sequent \lnot !A \lif \lnot !B$.
\item $\Sequent (!A \lif \lnot !A) \lif \lnot !A$.
\item $\Sequent \lnot(!A \lif !B) \lif \lnot !B$.
\item $!A \lif !C \Sequent \lnot (!A \land \lnot !C)$.
\item $!A \land \lnot !C \Sequent \lnot (!A \lif !C)$.
\item $!A \lor !B, \lnot !B \Sequent !A$.
\item $\lnot !A \lor \lnot !B \Sequent \lnot(!A \land !B)$.
\item $\Sequent (\lnot !A \land \lnot !B) \lif\lnot(!A \lor !B)$.
\item $\Sequent \lnot(!A \lor !B) \lif (\lnot !A \land \lnot !B)$.
\end{enumerate}
\end{prob}

\begin{prob}
નીચેના સિક્વન્ટોની !!{derivation}s આપો:
\begin{enumerate}
\item $\lnot(!A \lif !B) \Sequent !A$.
\item $\lnot(!A \land !B) \Sequent \lnot !A \lor \lnot !B$.
\item $!A \lif !B \Sequent \lnot !A \lor !B$.
\item $\Sequent \lnot \lnot !A \lif !A$.
\item $!A \lif !B, \lnot !A \lif !B \Sequent !B$.
\item $(!A \land !B) \lif !C \Sequent (!A \lif !C) \lor (!B \lif !C)$.
\item $(!A \lif !B) \lif !A \Sequent !A$.
\item $\Sequent (!A \lif !B) \lor (!B \lif !C)$.
\end{enumerate}
(આ બધામાં $\RightR{\Contraction}$ નિયમ જરૂરી છે.)
\end{prob}
\end{document}
